\documentclass[12pt,table,xcolor={dvipsnames}]{beamer}
\usetheme{Pittsburgh}
\usecolortheme{seagull}
%\usepackage[utf8]{inputenc}
\usepackage{fontspec}
\usepackage{amsmath}
\usepackage{listings}
\usepackage{multirow}
\usepackage{amsfonts}
\usepackage{amssymb}
\usepackage{graphicx}
\author{Design and Verification of Security Protocols and Security Ceremonies}
\title{Threat Modelling for Symbolic Evaluation}
%\setbeamercovered{transparent} 
\setbeamertemplate{navigation symbols}{} 
%\logo{\includegraphics[scale=0.015]{Brasao_UFSC.png}\includegraphics[scale=0.2]{brasao_PPGCC.jpg}} 
\institute{Programa de Pós-Graduacão em Ciências da Computacão \\ Dr. Jean Everson Martina} 
\date{\vspace{.2cm}March-June 2018} 
\subject{} 
\usebackgroundtemplate{\includegraphics[width=\paperwidth,
height=\paperheight]{../reusable_images/fundo_UFSC.png}}
\begin{document}

{
\usebackgroundtemplate{\includegraphics[width=\paperwidth,
height=\paperheight]{../reusable_images/fundo_capa.png}}
\begin{frame}
\titlepage
\includegraphics[scale=0.3]{../reusable_images/brasao_PPGCC.jpg}
\end{frame}
}

\frame{
	\frametitle{How to Start!}
\begin{itemize}
\item How can we tell if a cryptographic protocol or a security ceremony is secure? \pause
\item Phrased another way, how can we be sure that a given protocol/ceremony meets a given security goal?\pause
\item Whom are we defending against?
\end{itemize}
}

\frame{
	\frametitle{Notes about Security}
\begin{itemize}
\item The sense of security is directly related to the threats available within the environment;\pause
\item One would feel unsafe in a setting where others would feel secure;\pause
\item Is security a sensation?
\end{itemize}
}

\frame{
	\frametitle{Notes about Insecurity}
\begin{itemize}
\item As the security is related to the threats:\pause
\begin{itemize}
\item Assume no threats and you will be secure!\pause
\item Assume weaker threats and you still be secure;\pause
\item Are we able to foresee all the threats in a scenario?
\end{itemize}
\end{itemize}
}

\frame{
	\frametitle{Security and Insecurity comparison}
\begin{itemize}
\item We can have two different things, one that we consider secure and another that we don't;\pause
\item How do we compare two different things in terms of security?\pause
\item We need a baseline for comparison;\pause
\item The name of this baseline is a Threat Model. 
\end{itemize}
}

\frame{
	\frametitle{Changing the Baseline Matters}
\begin{itemize}
\item We can recall the NSPKP example:
\begin{itemize}
\item Needham and Schroeder assumed no internal could be an attacker;\pause
\item Everything worked for 15 years;\pause
\item Gavin Lowe assumed an internal could engage in the protocol to take advantage;\pause
\item The protocol is considered broken since.
\end{itemize} 
\end{itemize}
}

\frame{
	\frametitle{Importance of Threat Models}
\begin{block}{Lesson from NSSPK}
To claim security and not to be surprised in the future you need to clearly specify the threat model of your protocol or ceremony and it must be as close as the environment where the protocol or ceremony will run within. 
\end{block}
}

\frame{
	\frametitle{The Dolev-Yao Threat Model}
\begin{itemize}
\item D. Dolev and A. Yao. “On the security of public key protocols”. IEEE Transactions on Information Theory, 29(2):198-208. 1983;\pause
\item The original motivation for the paper was to verify public key protocols against active attackers with considerable power;\pause
\item The setting for the paper was the cold war times;\pause
\item Most mechanised formal methods for security analysis use some version of this model.
\end{itemize}
}

\frame{
	\frametitle{The Dolev-Yao Assumptions}
\begin{itemize}
\item The underlying public key system is 'perfectly secure':\pause
\begin{itemize}
\item One-way functions are unbreakable;\pause
\item Public directory is secure and cannot be tampered with;\pause
\item Everyone has access to all public keys;\pause
\item Only the peer knows his private key.
\end{itemize}
\end{itemize}
}

\frame{
	\frametitle{The Dolev-Yao Assumptions}
\begin{itemize}
\item Adversary has complete control over the entire network:\pause
\begin{itemize}
\item He acts as a legitimate user;\pause
\item He can obtain any message from any party;\pause
\item He can initiate the protocol with any party, and can be a receiver to any party in the network;\pause
\item Can read any message, decompose it into parts and re-assemble.
\end{itemize}
\end{itemize}
}

\frame{
	\frametitle{The Dolev-Yao Assumptions}
\begin{itemize}
\item Concurrent executions of the protocol can occur;\pause
\item The attacker cannot gain partial knowledge and perform statistical tests;\pause
\item The attacker can decrypt if and only if he knows the correct key;\pause
\item We assume that cryptographic functions have no special properties.
\end{itemize}
}


\frame{
	\frametitle{The Dolev-Yao Benefits}
\begin{itemize}
\item The model has the following features, which can be viewed as either benefits or restrictions depending on what you are trying to do:\pause
\begin{itemize}
\item It is simple to describe protocols in this model;\pause
\item Adversary has unlimited power, so although this is a conservative approach this may not be realistic;\pause
\item Protocols have a 'black-box' nature, which means that linking individual protocols with others is extremely difficult. Dolev-Yao helps on that.
\end{itemize}
\end{itemize}
}

\frame{
	\frametitle{The Dolev-Yao Features}
\begin{itemize}
\item Secrecy properties: Alice is given a message M as input, and starts exchanging messages with Bob, with the goal of sending message M to Bob in a secret way;\pause
\item The security property is that an adversary cannot recover M, even if actively interfering with the protocol;\pause
\item No other security properties are considered.
\end{itemize}
}

\frame{
	\frametitle{The Dolev-Yao Features}
\begin{itemize}
\item Stateless parties: The main limitation on the honest parties is that they are stateless;\pause 
\item The messages transmitted by a party at every step of the protocol are a function of their initial knowledge and the message they just received;\pause
\item In particular, parties cannot use information collected from previous messages not addressed to them; \pause
\item For this reason, these protocols have been named "ping-pong" protocols;\pause 
\item We observe that the stateless restriction is only put on the honest parties, and the adversary can maintain state, record communications, and store values that are subsequently used in the construction of messages;
\end{itemize}
}

\frame{
	\frametitle{The Dolev-Yao Features}
\begin{itemize}
\item Concurrent execution: The adversary can start an arbitrary number of protocol executions, involving different sets of parties, where each player can participate in several concurrent executions;\pause
\item In this respect, the model considered here is more general than the computational model considered at the time, which focused on single protocol execution;\pause
\item The computational cryptography community started addressing the important issue of concurrency only in the 90's.
\end{itemize}
}

\frame{
	\frametitle{The Dolev-Yao Features}
\begin{itemize}
\item In the Dolev-Yao model, there are two kinds of active parties: honest participants
and the adversary; \pause
\item The honest participants follow the steps of the protocol without deviation;\pause
\item The attacker does not follow the rules;\pause
\item The peers do not share long term secrets, even the attacker keeps things for himself.
\end{itemize}
}

\frame{
	\frametitle{ Dolev-Yao in Practice}
\begin{itemize}
\item All the peers start with some knowledge set, including all public information, peers names and their own random numbers to be used as nonces;\pause
\item If a message is cast in the network, it go to the destination's knowledge set and the attackers knowledge set (Eavesdrop);
\end{itemize}
}


\frame{
	\frametitle{ Dolev-Yao in Practice}
\begin{itemize}
\item From the attacker knowledge set:
\begin{itemize}
\item He can decrypt whatever he has the key to do so and re-encrypt whatever he knows with the keys he has at hand (Encryption);\pause
\item He can break down messages to atomic components (Atomic Breakdown);\pause
\item He can re-arrange all the components he knows in all possible ways to form new messages (Fabricate);
\end{itemize}
\end{itemize}
}

\frame{
	\frametitle{ Dolev-Yao in Practice}
\begin{itemize}
\item He can modify messages in transit (Modify);\pause
\item He can replay any message he collects from traffic (Replay);\pause
\item He can prevent the delivery of any message (Block);\pause
\item He can engage legitimately with other peers on the protocol (Initiate);\pause
\item He can re-arrange the order of the messages in traffic (Re-order);\pause
\item He can mimic the identity of any peer in the network (Spoof). 
\end{itemize}
}

\frame{
	\frametitle{Dolev-Yao Implementation}
\begin{itemize}
\item It is usually logically implemented;\pause
\item It uses free variables to allow for the powers to happen;\pause
\item Although the amount of knowledge the attacker can have is potentially unbound, it is finite;\pause
\item Some techniques create other significant limits to the Dolev-Yao Threat Model.
\end{itemize}
}

\frame{
	\frametitle{Discussion}
\begin{itemize}
\item Is Dolev-Yao enough to put a security protocols against and claim it secure?\pause
\item Can you foresee some capability missing?\pause
\item Is Dolev-Yao sufficient to compare two security protocols?\pause
\item What can go wrong when we compare two security protocols that use Dolev-Yao as Threat Model?
\end{itemize}
}

\frame{
	\frametitle{Dolev-Yao Considerations}
\begin{itemize}
\item DY can be considered the standard threat model to study security protocols;\pause. \item The DY attacker controls the entire network although he cannot perform cryptanalysis;\pause
\item The DY model has remarkably favoured the discovery of significant protocol flaws,
but the attacker has significantly changed today;\pause
\item To become an attacker has never been so easy.
\end{itemize}
}

\frame{
	\frametitle{A Family of Variations for Dolev-Yao}
\begin{itemize}
\item B.U.G;\pause
\item The Rational Attacker;\pause
\item The General Attacker;\pause
\item Multi Attacker;\pause
\item Distributed Attacker.
\end{itemize}
}

\frame{
	\frametitle{The B.U.G Threat Model}
\begin{itemize}
\item B.U.G. dates back to 2002;\pause
\item The name is a permuted acronym for the “Good”, the “Bad” and the “Ugly”;\pause
\item This model attempts stricter adherence to reality by partitioning the participants into three groups;\pause
\item The Good principals would follow the protocol;\pause
\item The Bad would in addition try to subvert it;\pause
\item The Ugly would be ready to either behaviour.
\end{itemize}
}

\frame{
	\frametitle{The B.U.G Threat Model Insights}
\begin{itemize}
\item A principal may change role and decide to attempt illegal exploitation of a protocol although he has always conformed to it so far;\pause
\item It changes the idea of a single attacker since anyone could attack;\pause
\item The principle behind this threat model is that the attackers do not share long term secrets.
\end{itemize}
}

\frame{
	\frametitle{The B.U.G Threat Model Issues}
\begin{itemize}
\item It is unclear on how dynamic updates should be on the behaviour:\pause
\begin{itemize}
\item One could change behaviour after every single message;\pause
\item Either sent, received or cast;\pause
\end{itemize}
\item This complicates a lot the mechanisations of the attacker, because all behaviour is possible.
\end{itemize}
}

\frame{
	\frametitle{The Rational Attacker Threat Model}
\begin{itemize}
\item BUG appeared overly detailed, and was simplified as The Rational Attacker Threat
Model;\pause
\item It was conceived in 2008;\pause
\item The Rational Attacker let any principal make cost/benefit decisions at any time
to either behave according to the protocol or not;\pause
\item Analysing a protocol under the Rational Attacker requires specifying each principal’s
cost and benefit functions.
\end{itemize}
}

\frame{
	\frametitle{The Rational Attacker Insights}
\begin{itemize}
\item The Rational Attacker  seems out of reach for the current mechanised approaches, especially for bound verification techniques;\pause
\item Although complex to mechanise, the Rational Attacker is more realistic than B.U.G.;\pause
\item In the wild, it is common to the attacker to make cost/benefit analysis when to engage or not;\pause
\item The Rational Attacker bring all game theory into the protocols' scenarios.
\end{itemize}
}

\frame{
	\frametitle{The Rational Attacker Issues}
\begin{itemize}
\item The Rational Attacker is not clear whether the cost/benefit function is fixed or variable;\pause
\item One would argue that the objectives of the attacker are not static and that it would change depending on the gains made so far;\pause
\item Mechanisation is not only and issue of representativeness of the formal verification technique, but an entangled problem.
\end{itemize}
}

\frame{
	\frametitle{The General Attacker Threat Model}
\begin{itemize}
\item  The General Attacker abstracts away the actual cost/benefit analysis in a simplified model;\pause
\item Any principal may behave as a Dolev-Yao attacker;\pause
\item The change of perspective in RA or in GA with respect to DY is clear: principals
do not collude for a common aim but, rather, each of them acts for his own personal
sake;\pause
\item By contrast, a pair of colluding DY attackers is equivalent to a single DY
attacker in terms of generated attacks;\pause
\begin{itemize}
\item This is confirmed by a formal proof.
\end{itemize}
\end{itemize}
}

\frame{
	\frametitle{The General Attacker Insights}
\begin{itemize}
\item  Endowing each principal with the entire potential of a DY attacker signifies that he may send any of the messages he can form to anyone;\pause
\item Such messages include both the legal ones, conforming to the protocol in use, and the illegal, forged ones, which he can build from the analysis of the traffic though without cryptanalysis.
\end{itemize}
}

\frame{
	\frametitle{The General Attacker Issues}
\begin{itemize}
\item  It was suggested that the General Attacker is similar to Dolev-Yao provided that all principals reveal their secrets to the attacker (Augmented Dolev-Yao); \pause
\item They appear equivalent: \pause
\begin{itemize}
\item Any illegal message that a principal may send in General Attacker may be sent by the
single augmented Dolev-Yao attacker;\pause
\item This happens because he knows everyone’s secrets.
\end{itemize}  
\end{itemize}
}

\frame{
	\frametitle{The General Attacker versus Augmented Dolev-Yao}
\begin{itemize}
\item Augmented Dolev-Yao entangles the interpretation of attacks where principals attack each other;\pause
\item The single attacker will always be the originator of any attack, complicating the  identification of the real perpetrator;\pause
\item For attacks against the attacker, the model will feature the attacker attacking himself, thus stretching the interpretation of the victim to an extreme;\pause 
\item Perpetrator and victim are naturally expressed in GA because its gist is exactly to reflect modern everyone-for-themselves scenarios.
\end{itemize}
}

\frame{
	\frametitle{The Multi Attacker Threat Model}
\begin{itemize}
\item In the Multi Attacker each principal may behave as a Dolev-Yao attacker but will never reveal his long-term secrets;\pause
\item It was conceived in 2011;\pause
\item Multi Attacker  can be seen as a refinement of General Attacker with some rationality that avoids the trivial impersonation attacks;\pause
\item It helps to understand some new types of attacks.
\end{itemize}
}

\frame{
	\frametitle{The Multi Attacker Insights}
\begin{itemize}
\item Analysing protocols under the Multi Attacker threat model yields unknown scenarios of retaliation or anticipation;\pause
\item If an attack can be retaliated under Multi Attacker, such a scenario will not occur under Rational Attacker because the cost of attacking clearly overdoes its benefit, and hence the attacker will not attack in the first place;\pause
\item This changes the game of how a powerful attacker would attack, because retaliation may let the attacker vulnerable.
\end{itemize}
}

\frame{
	\frametitle{The Multi Attacker Issues}
\begin{itemize}
\item The Multi Attacker do not use its full capabilities to derive partial information;\pause
\item By being powerful and knowing what is going on, he could anticipate what other Multi Attacker have on their knowledge set;\pause
\item This is not encoded on the attacker;\pause
\item This would make competition between the attacker fiercer.
\end{itemize}
}

\frame{
	\frametitle{The Distributed Attacker Threat Model}
\begin{itemize}
\item The Distributed Attacker is an evolution of the Multi Attacker where the principals can have different powers;\pause
\item It was conceived in 2015;\pause
\item It was tailored for a layered strategy for security ceremonies;\pause
\item It is reasonable to use in protocol verification because the real capabilities of the multi-attacker we have are not clear and eventually will not be the same.
\end{itemize}
}

\frame{
	\frametitle{The Distributed Attacker Insights}
\begin{itemize}
\item It is based on Martina-Carlos ideas of breaking down the power of the Dolev-Yao attacker;\pause
\item Acknowledging that each multi-attacker has different powers it a good strategy that can show us the competition between the attacker for the target;\pause
\item It was shown that we can mechanise such attacker using First-Order Logics;\pause
\item No issues were brought so far due to its freshness within the protocol and ceremony verification communities.
\end{itemize}
}
 
\frame{
	\frametitle{Discussion}
\begin{itemize}
\item Can you identify a protocol where using one of this evolutions of Dolev-Yao can bring insights of new attacks?\pause
\item Which one would make more sense today?\pause
\item Choosing one Threat Model to work invalidate the others?
\end{itemize}
}

{
\usebackgroundtemplate{\includegraphics[width=\paperwidth,
height=\paperheight]{../reusable_images/fundo_capa.png}}
\begin{frame}

{\LARGE Questions????}

\end{frame}
}

{
\usebackgroundtemplate{\includegraphics[width=\paperwidth,
height=\paperheight]{../reusable_images/fundo_capa.png}}
\begin{frame}
\includegraphics[scale=0.8]{../reusable_images/cc_logo_arge.png}\hspace{0.5cm} 
\includegraphics[scale=0.95]{../reusable_images/by.png}

\vspace{1cm}
This work is licensed under the Creative Commons Attribution 4.0 International License. To view a copy of this license, visit http://creativecommons.org/licenses/by/4.0/.
\end{frame}
}

\end{document}